\clearpage
\fancyhead[L]{\small GAUSS AN VON LINDENAU.}
\hypertarget{native-301-lower}{}
\begin{center}GAUSS AN VON LINDENAU.
\end{center}

\begin{center}
\begin{tabular}{c}
\hline
Monatliche Correspondenz. Band XVII. S. 182--188. 1808 Februar.\\
\hline
\end{tabular}
\end{center}

\begin{flushright}
Göttingen, am 25. Januar 1808.
\end{flushright}

Bitte entschuldigen Sie, dass ich Ihnen die erhoffte Reihe guter hiesiger Kometenbeobachtungen wegen des beispiellos schlechten Wetters nicht mitschicken kann. In der ganzen Zeit, die ich hier bin (seit dem 21. November), gab es nur 3 Abende, an denen ich dem ungünstigen Himmel einige spärliche Beobachtungen abgewinnen konnte. Ich gebe sie hier zusammen mit den Abweichungen von meinen Elementen wieder. [Es sind die Beobachtungen vom 17. und 29. Dezember und vom 4. Januar, die weiter unten auf Seite 298 im Abdruck des Aufsatzes vom 25. Februar 1808 aus den Göttingischen gelehrten Anzeigen stehen.]

\par\indent Auf die Beobachtungen vom 29. Dezember kann man sich weniger verlassen, weil das dafür verwendete Instrument schlecht montiert ist. Die beiden ersten Beobachtungen vom 4. Januar sind nur einzelne Bestimmungen zwischen den Wolken. Für sich genommen halte ich sie aber für gut. Die letzte stammt von meinem geschätzten Kollegen, Prof. \textsc{Harding}, und ist der Mittelwert aus einigen gut übereinstimmenden Vergleichen. Im Übrigen können die jüngsten Beobachtungen des Kometen keine große Genauigkeit erreichen, da seine geringe Lichtstärke und sein kaum erkennbarer Kern keine so genaue Beobachtung der Ein- und Austritte mehr zulassen. Im Dezemberheft der M.\,C. [S. 292 dieses Bandes] ist bei meinen Elementen übrigens die Bahnneigung vergessen worden. Sie beträgt $63^\circ\;12'\;35''6$. Das sehe ich in dem Exemplar von Hofrat \textsc{Blumenbach}, das mir soeben gezeigt wurde.

\par\indent Für die Mailänder Beobachtungen der \VestaAltarThreeOhOne{}, die Sie mir freundlicherweise mitgeteilt haben, bin ich Ihnen sehr verbunden. Mein vortrefflicher Freund Herr \textsc{Bessel} war so freundlich, die letzten Beobachtungen unter gebührender Berücksichtigung der Refraktion so genau wie möglich zu reduzieren. Hier sind die Ergebnisse dieser Reduktionen:


\begin{center}\fontsize{9}{11}\selectfont\setlength{\tabcolsep}{2pt}\renewcommand{\arraystretch}{1.15}\begin{tabular}{lrrrrr}
& & & & \multicolumn{2}{c}{Differenzen}\\
1807 & Mittl. Zeit & Rektasz. & Deklination & Rektasz. & Dekl.\\\hline
Sept. 5 & $7^{\text{u}}\,21^{\text{m}}\,13^{\text{s}}0$ & $211^\circ\;29'\;4''5$ & $8^\circ\;2'\;49''9$ südl. & $-46''6$ & $-53''0$\\
8 & $7\;14\;47.5$ & $212\;49\;55.4$ & $8\;38\;18.1$ & $-54.0$ & $-59.3$\\
8 & $7\;29\;44.8$ & $212\;50\;11.5$ & $8\;38\;16.0$ & $-53.2$ & $-49.9$\\
10 & $7\;13\;14.9$ & $213\;44\;31.7$ & $9\;1\;45.8$ & $-65.5$ & $-61.4$\\
21 & $7\;4\;10.8$ & $218\;51\;35.2$ & $11\;7\;42.5$ & $-82.4$ & $-64.5$\\
26 & $6\;53\;51.8$ & $221\;15\;21.1$ & $12\;2\;59.3$ & $-99.5$ & $-66.9$\\
\end{tabular}\end{center}

\par\indent Prof. \textsc{Harding} war so freundlich, diese Beobachtungen mit meinen III. Elementen zu vergleichen. Daraus ergaben sich die folgenden Differenzen. [In diesem Abdruck stehen sie oben neben den Beobachtungen.]

\par\indent I\SourceLeaf{302}ch halte es für überflüssig, die III. Elemente auf dieser Grundlage jetzt noch einmal zu verbessern. Die Korrekturen wären nur sehr klein, und die Ephemeride im Septemberheft der M.\,C. wird zum Auffinden genügen.

\par\indent Hier sind die neuen Elemente der \emph{Pallas}, nach denen die Ephemeride im Dezemberheft der M.\,C. berechnet wurde:

\begin{center}
X$^{\text{te}}$ Elemente
\end{center}


\begin{center}\fontsize{10}{12}\selectfont\setlength{\tabcolsep}{2pt}\renewcommand{\arraystretch}{1.15}\begin{tabular}{cr}
Epoche & \shortstack{Meridian von\\Göttingen}\\\hline
1802 & $143^\circ\;33'\;58''7$\\
1803 & $221\;39\;26.9$\\
1804 & $299\;57\;45.4$\\
1805 & $18\;3\;13.6$\\
1806 & $96\;8\;41.8$\\
1807 & $174\;14\;10.0$\\
1808 & $252\;32\;28.5$\\
\end{tabular}\end{center}

\begin{center}\fontsize{10}{12}\selectfont\setlength{\tabcolsep}{2pt}\renewcommand{\arraystretch}{1.15}\begin{tabular}{lr}
Tägliche tropische Bewegung .\;.\;.\;.\;.\; & $770^{.}2143$\\
Logarithmus der großen Halbachse .\;.\; & $0.4423149$\\
Exzentrizität .\;.\;.\;.\;.\;.\;.\;.\;.\;.\;.\; & $0.2450198$\\
Perihel 1803 .\;.\;.\;.\;.\;.\;.\; & $121^\circ\;3'\;11''4$\\
$\AscendingNode$, 1803 .\;.\;.\;.\;.\;.\;.\;.\;.\;.\; & $172\;28\;56.9$\\
Neigung der Bahn .\;.\;.\;.\;.\;.\;.\; & $34\;37\;41.0$\\
\end{tabular}\end{center}

\par\indent Diese Elemente stützen sich ausschließlich auf die vier Oppositionen von 1803, 1804, 1805 und 1807. Die Beobachtungen von 1802 wurden überhaupt nicht berücksichtigt, da sie sich für eine genaue Bestimmung der Opposition nicht eignen. Alle Längen jener Oppositionen werden genau wiedergegeben,, die Breiten lassen sich zwar nicht genau in Einklang bringen, doch erreichen die Differenzen keine $\frac12$ Minute. Im Übrigen wird Pallas 1808 schwieriger zu beobachten sein als je zuvor, da sie gerade im Aphel steht.

\Needspace{19\baselineskip}
\par\indent Vor einiger Zeit habe ich auch die Elemente der \emph{Juno} verbessert. Die Ergebnisse und die daraus berechnete Ephemeride für 1808 finden Sie in der Beilage. Ich möchte nur noch erwähnen, dass Prof. \textsc{Harding} so freundlich war, die neuen VII$^{\text{ten}}$ Elemente sorgfältig mit den Beobachtungen von Herrn \textsc{Bessel} zu vergleichen, wobei er die Parallaxe außer Acht ließ.


\begin{center}\fontsize{9}{11}\selectfont\setlength{\tabcolsep}{2pt}\renewcommand{\arraystretch}{1.15}\begin{tabular}{lrrrrr}
1807 & \shortstack{Mittlere Zeit\\Lilienthal} & \shortstack{Berechnete\\Rektasz.} & \shortstack{Untersch.\\R--B} & \shortstack{Berechn. südl.\\Deklination} & \shortstack{Untersch.\\R--B}\\\hline
Mai 5 & $12^{\text{u}}\,46^{\text{m}}\,4^{\text{s}}$ & $240^\circ\;5'\;11''4$ & $+26''6$ & $3^\circ\;44'\;27''1$ & $+7''8$\\
7 & $12\;25\;33$ & $239\;42\;31.0$ & $+3.6$ & $3\;34\;18.0$ & $+17.6$\\
8 & $11\;30\;29$ & $239\;31\;22.3$ & $+5.6$ & $3\;29\;29.4$ & $+0.1$\\
19 & $12\;8\;26$ & $237\;17\;26.9$ & $-5.6$ & $2\;40\;34.5$ & $-3.8$\\
23 & $11\;51\;8$ & $236\;28\;5.2$ & $-9.5$ & $2\;26\;12.1$ & $+5.9$\\
24 & $11\;14\;57$ & $236\;16\;7.4$ & $-5.2$ & $2\;22\;59.9$ & $-0.2$\\
25 & $11\;54\;26$ & $236\;3\;31.0$ & $-2.6$ & $2\;19\;44.7$ & $-2.3$\\
26 & $11\;50\;58$ & $235\;51\;11.4$ & $-5.7$ & $2\;16\;41.9$ & $-16.0$\\
\end{tabular}\end{center}

\par\indent In der Ephemeride für die $\PallasSymbol$ habe ich einen Druckfehler zu berichtigen. Die letzte Spalte gibt nicht die Abstände der Pallas von der Erde an, sondern deren Logarithmen.

\par\indent Mein Werk über die Bestimmung der Bahnen der Himmelskörper ist jetzt im Druck und wird hoffentlich bald nach Ostern erscheinen können.

\par\indent Ich schicke Ihnen auch kleine Aberrations- und Nutationstafeln mit. Ich hatte sie schon vor vier Jahren berechnet und mehreren Freunden mitgeteilt, die sie praktisch fanden und sich an ihren Gebrauch gewöhnt haben. Mir scheint, dass man mit ihnen von allen allgemeinen Tafeln am schnellsten und kürzesten rechnet. Dass man dafür Sinus- und Logarithmentafeln braucht, fällt nicht ins Gewicht, denn diese hat ohnehin jeder Astronom stets auf dem Schreibtisch. Das Exemplar die\SourceLeaf{303}ser Tafeln, das ich Ihnen hier schicke, ist nach den neuesten Zahlenangaben von Herrn \textsc{von Zach} ganz neu berechnet worden. [Seite 123..128 dieses Bandes.]

\par\indent [Der geozentrische Lauf der Juno ist für Mitternacht in Göttingen vom 16. April bis zum 28. Dezember 1808 in Abständen von vier Tagen angegeben, die Rektaszension und Deklination auf Minuten genau und der Logarithmus des Abstands von der Erde auf vier Dezimalstellen.]

